\newpage
\chapter{PENDAHULUAN} \label{Bab I}

\section{Latar Belakang} \label{I.Latar Belakang}
Dalam beberapa dekade terakhir, kemajuan dalam pemrosesan suara dan pengenalan ucapan (\textit{Automatic Speech Recognition} / ASR) telah menjadi salah satu bidang riset yang sangat penting, terutama karena potensinya besar di banyak aspek kehidupan manusia, misalnya asisten suara (\textit{Siri, Google Assistant}), sistem transkripsi otomatis, aksesibilitas bagi penyandang disabilitas, layanan publik, pendidikan, serta interaksi manusia-komputer. Model-model berbasis pembelajaran mendalam (\textit{deep learning}) yang menggunakan arsitektur \textit{Transformer}, atau model "\textit{end-to-end}", semakin banyak digunakan untuk menggantikan pendekatan tradisional seperti \textit{Hidden Markov Model dan Gaussian Mixture Model}. \par

Namun, performa ASR sangat bergantung pada ketersediaan data berbicara (\textit{speech}) dan transkripsi yang berkualitas dan representatif dari bahasa atau dialek target. Untuk bahasa mayor seperti Inggris, Mandarin, Spanyol, data sangat melimpah; sementara untuk bahasa lokal atau dialek (termasuk bahasa daerah di Indonesia) data seringkali sangat terbatas (\textit{low-resource}). Model yang telah dilatih pada data multibahasa besar seperti Whisper kadang masih kurang optimal ketika menghadapi dialek atau bahasa yang kurang terwakili dalam data latihannya. \par

OpenAI memperkenalkan model Whisper sebagai model \textit{speech-to-text} multibahasa dengan pelatihan lemah-supervisi skala besar (\textit{Large-Scale Weak Supervision}) (sekitar 680.000 jam data audio transkrip) untuk generalisasi ke berbagai kondisi dan bahasa \cite{radford2023whisper}. Meskipun demikian, model dasar ("\textit{base Whisper}") seringkali tidak menghasilkan akurasi memadai ketika dihadapkan pada bahasa atau dialek yang sangat spesifik atau dengan variasi fonetik tinggi yang belum banyak direpresentasikan. Oleh karena itu, teknik \textit{fine-tuning} terhadap model dasar menjadi sangat penting agar model "teradaptasi" ke karakteristik bahasa lokal tersebut. \par

Teknik \textit{fine-tuning} Whisper telah dieksplorasi dalam sejumlah studi, terutama untuk bahasa \textit{low-resource}. Misalnya, penelitian \textit{Exploration of Whisper fine-tuning strategies for low-resource ASR} menunjukkan bahwa berbagai strategi \textit{fine-tuning} dapat secara signifikan meningkatkan performa pengenalan suara untuk bahasa-bahasa dengan data terbatas \cite{liu2024exploration}. Ada pula penelitian "\textit{Fine-tuning Whisper on Low-Resource Languages for Real-World Applications}" yang memperkenalkan metode transformasi dataset kalimat-ke-audio panjang untuk menjaga kemampuan segmentasi dan bagus bagi aplikasi nyata \cite{timmel2024finetune}. \par

Studi lain mengusulkan metode \textit{fine-tuning} efisien, seperti \textit{LoRA-Whisper}, yang bertujuan agar adaptasi ke bahasa baru bisa dilakukan tanpa menurunkan performa bahasa lain dan dengan \textit{overhead} parameter yang lebih kecil \cite{song2024lorawhisper}. Juga ada penelitian \textit{Multistage Fine-tuning Strategies for Automatic Speech Recognition in Low-resource Languages} yang menggunakan strategi \textit{multistage} (bertahap) untuk mengadaptasi model melalui bahasa yang punya kemiripan linguistik \cite{pillai2024multistage}. \par

Begitu juga, studi aplikasi domain khusus (misalnya dalam konteks medis) menekankan bahwa \textit{fine-tuning} model Whisper pada domain spesifik (dengan kosakata dan pola suara khusus) mampu meningkatkan akurasi dibandingkan model standar \cite{banerjee2024unitedmedasr}. Dengan demikian, adaptasi model ASR (melalui \textit{fine-tuning}) telah menjadi pendekatan penting untuk menjembatani kesenjangan antara model umum dan kebutuhan lokal atau domain khusus. \par

Di Indonesia, keragaman bahasa dan dialek sangat tinggi, dan banyak komunitas menggunakan bahasa daerah dalam kehidupan sehari-hari. Bahasa Minangkabau adalah salah satu bahasa daerah yang memiliki nilai budaya, sosial, dan historis tinggi, khususnya di wilayah Sumatera Barat dan daerah-daerah diaspora Minangkabau. Namun, dalam literatur riset ASR Indonesia, fokus utama lebih banyak pada bahasa Indonesia (standar) atau beberapa dialek besar (Sunda, Jawa, dsb.), sedangkan kontribusi penelitian terhadap bahasa Minangkabau relatif sedikit. \par

Karena bahasa Minangkabau memiliki fonetik, kosakata, dan intonasi yang khas, transkripsi otomatis dengan model umum (bahasa Indonesia atau multibahasa) cenderung menghasilkan kesalahan (\textit{error}) yang cukup tinggi. Hal ini dapat disebabkan oleh kekurangan data latih (audio + transkripsi) dalam bahasa Minangkabau, serta ketidakmampuan model awal menangani variasi dialek, aksen lokal, dan fenomena fonologi spesifik. Oleh karena itu, tugas untuk mengembangkan sistem ASR yang mampu mengenali ucapan bahasa Minangkabau dengan baik memerlukan adaptasi model melalui \textit{fine-tuning}, agar model "belajar" karakteristik suara, kosakata, dan pola ucapan lokal tersebut. \par

Urgensi pengembangan sistem ASR untuk bahasa Minangkabau semakin meningkat seiring dengan kebutuhan nyata di berbagai sektor. Di bidang \textbf{pelestarian budaya}, terdapat banyak arsip audio sastra lisan Minangkabau (seperti \textit{kaba}, \textit{pantun}, dan cerita rakyat) yang tersimpan dalam format analog atau rekaman digital tanpa transkripsi, sehingga sulit diakses dan dianalisis oleh generasi muda maupun peneliti. Sistem ASR yang akurat dapat mempercepat \textbf{digitalisasi dan dokumentasi} warisan budaya ini. Di sektor \textbf{media dan konten digital}, para \textit{content creator} lokal yang memproduksi konten video atau podcast berbahasa Minangkabau memerlukan alat transkripsi otomatis untuk membuat subtitle atau transkrip guna meningkatkan aksesibilitas dan jangkauan audiens. Selain itu, di bidang \textbf{pariwisata dan layanan publik}, sistem ASR dapat diintegrasikan ke dalam aplikasi pemandu wisata interaktif atau layanan informasi berbasis suara untuk wisatawan yang ingin memahami budaya Minangkabau. Dengan demikian, keberadaan sistem ASR bahasa Minangkabau bukan hanya bersifat \textit{"nice to have"} untuk riset akademis, tetapi merupakan kebutuhan \textit{essential} yang memiliki dampak praktis langsung terhadap pelestarian budaya, ekonomi kreatif, dan aksesibilitas informasi bagi masyarakat Minangkabau. \par

\section{Rumusan Masalah} \label{I.Rumusan Masalah}
Berdasarkan latar belakang yang telah dijelaskan, maka permasalahan utama dalam penelitian ini dapat dirumuskan sebagai berikut:

\begin{enumerate}
    \item Bagaimana mengadaptasi (\textit{fine-tuning}) model \textbf{OpenAI Whisper Small} agar mampu mengenali dan mentranskripsi ucapan dalam \textbf{bahasa Minangkabau} secara akurat, mengingat keterbatasan data suara dan teks (\textit{low-resource}) yang tersedia?
    
    \item Bagaimana proses perancangan dan implementasi sistem \textbf{Automatic Speech Recognition (ASR)} berbasis hasil \textit{fine-tuning} tersebut dengan menggunakan bahasa pemrograman \textbf{Python} dan pustaka \textit{machine learning} seperti \textbf{PyTorch} dan \textbf{Hugging Face Transformers}?
    
    \item Bagaimana pengaruh penerapan teknik \textit{fine-tuning} tertentu (misalnya \textit{full fine-tuning} atau \textit{parameter-efficient fine-tuning} seperti \textbf{LoRA}) terhadap \textbf{kinerja model Whisper Small} dalam mengenali ucapan bahasa Minangkabau, terutama dalam konteks keterbatasan sumber daya komputasi dan ukuran dataset yang terbatas pada bahasa \textit{low-resource}?
    
    \item Bagaimana cara mengukur dan mengevaluasi \textbf{performa model hasil fine-tuning} menggunakan metrik kuantitatif seperti \textit{Word Error Rate (WER)} dan \textit{Character Error Rate (CER)}, serta melakukan analisis kualitatif terhadap kesalahan transkripsi dengan mengkategorikan jenis kesalahan (fonetik, leksikal, dialek-spesifik, dan kontekstual) pada beragam variasi penutur dan kondisi rekaman?
    
    \item Bagaimana hasil implementasi dan pengujian sistem ASR yang dikembangkan dapat memberikan \textbf{kontribusi nyata terhadap pelestarian bahasa daerah} dan pengembangan teknologi pengenalan suara untuk bahasa lokal di Indonesia?
\end{enumerate}

\section{Tujuan Penelitian} \label{I.Tujuan}
Berdasarkan rumusan masalah yang telah dikemukakan, penelitian ini memiliki beberapa tujuan sebagai berikut:

\begin{enumerate}
    \item \textbf{Mengadaptasi model OpenAI Whisper Small} melalui proses \textit{fine-tuning} agar mampu mengenali dan mentranskripsi ucapan dalam \textbf{bahasa Minangkabau} dengan tingkat akurasi yang lebih baik dibandingkan model dasar (\textit{pretrained model}).
    
    \item \textbf{Merancang dan mengimplementasikan sistem Automatic Speech Recognition (ASR)} berbasis hasil \textit{fine-tuning} model Whisper menggunakan bahasa pemrograman \textbf{Python} dan pustaka \textit{deep learning} seperti \textbf{PyTorch} serta \textbf{Hugging Face Transformers}.
    
    \item \textbf{Menerapkan dan membandingkan teknik fine-tuning yang berbeda}, seperti \textit{full fine-tuning} dan \textit{parameter-efficient fine-tuning} (misalnya \textbf{LoRA-Whisper}), untuk menentukan metode yang paling efektif dan efisien dalam konteks bahasa dengan sumber daya terbatas (\textit{low-resource language}), dengan mempertimbangkan trade-off antara akurasi, waktu pelatihan, dan kebutuhan memori komputasi.
    
    \item \textbf{Mengevaluasi performa model hasil fine-tuning} menggunakan metrik kuantitatif seperti \textit{Word Error Rate (WER)} dan \textit{Character Error Rate (CER)}, serta melakukan analisis kualitatif dengan mengkategorikan kesalahan transkripsi berdasarkan tipe kesalahan (fonetik, leksikal, dialek-spesifik, dan kontekstual) untuk mengidentifikasi pola kesalahan dominan dan memberikan rekomendasi perbaikan model.
    
    \item \textbf{Menghasilkan prototipe sistem ASR bahasa Minangkabau yang fungsional dan terverifikasi}, serta memberikan kontribusi terhadap upaya pelestarian bahasa daerah melalui penerapan teknologi kecerdasan buatan di bidang pengenalan suara multibahasa.
\end{enumerate}

\section{Batasan Masalah} \label{I.Batasan}
Agar penelitian ini memiliki ruang lingkup yang terarah dan dapat diselesaikan secara realistis dalam waktu yang terbatas, maka ditetapkan beberapa batasan masalah sebagai berikut:

\begin{enumerate}
    \item Penelitian ini hanya berfokus pada \textbf{implementasi \textit{fine-tuning} model OpenAI Whisper Small} untuk pengenalan ucapan (\textit{Automatic Speech Recognition}) pada \textbf{bahasa Minangkabau}.
    
    \item Dataset yang digunakan dalam penelitian ini bersumber dari \textbf{Librivox Indonesia}, sebuah korpus audio publik yang menyediakan rekaman suara berbahasa Minangkabau beserta transkripsi teksnya. Dataset ini mencakup rekaman dari berbagai penutur dengan karakteristik demografi yang beragam (usia, jenis kelamin, dan variasi dialek Minangkabau seperti dialek Agam, Tanah Datar, Pesisir, dan 50 Kota). Target volume data yang digunakan untuk pelatihan adalah sekitar \textbf{10-20 jam audio} dengan kualitas rekaman yang bervariasi untuk mensimulasikan kondisi nyata. Penelitian ini tidak mencakup proses pembuatan atau pengumpulan dataset baru dalam skala besar, melainkan fokus pada pemanfaatan dan preprocessing data yang sudah tersedia.
    
    \item Model yang diujikan dibatasi pada \textbf{varian Whisper Small} dan teknik \textit{fine-tuning} tertentu seperti \textit{full fine-tuning} dan \textit{parameter-efficient fine-tuning} (misalnya LoRA), tanpa membandingkan seluruh varian Whisper lainnya (Tiny, Base, Medium, Large).

    \item Evaluasi performa model dilakukan menggunakan metrik kuantitatif \textbf{Word Error Rate (WER)} dan \textbf{Character Error Rate (CER)}. Selain itu, dilakukan analisis kualitatif terhadap kesalahan transkripsi dengan mengkategorikan jenis kesalahan ke dalam beberapa tipe, yaitu: (1) \textbf{kesalahan fonetik} (misalnya kesalahan pengenalan fonem vokal atau konsonan khas Minangkabau), (2) \textbf{kesalahan leksikal} (misalnya kesalahan pada kata serapan dari bahasa Indonesia atau bahasa asing yang diucapkan dengan aksen Minangkabau), (3) \textbf{kesalahan terkait dialek} (misalnya perbedaan pengucapan antar dialek Agam, Tanah Datar, Pesisir, dan 50 Kota), serta (4) \textbf{kesalahan kontekstual} (kesalahan pada kata yang bergantung pada konteks kalimat). Analisis kualitatif ini akan dilakukan secara manual terhadap sampel hasil transkripsi untuk mengidentifikasi pola kesalahan yang dominan dan memberikan rekomendasi perbaikan model. Evaluasi subjektif berbasis pengguna akhir (\textit{user study}) tidak termasuk dalam ruang lingkup penelitian ini.

    \item Implementasi sistem ASR dilakukan menggunakan bahasa pemrograman \textbf{Python} dengan pustaka \textbf{PyTorch} dan \textbf{Hugging Face Transformers}, serta dijalankan dalam lingkungan komputasi lokal atau GPU yang tersedia. Integrasi ke aplikasi lain (misalnya mobile atau web) tidak menjadi fokus utama penelitian ini.
    
    \item Penelitian ini tidak membahas aspek keamanan data, privasi pengguna, maupun optimasi energi selama pelatihan model, melainkan berfokus pada \textbf{adaptasi model dan analisis performa hasil fine-tuning}.
\end{enumerate}


\section{Manfaat Penelitian} \label{I.Manfaat}
Penelitian ini diharapkan dapat memberikan manfaat baik dari segi akademik maupun praktis, yaitu sebagai berikut:

\begin{enumerate}
    \item \textbf{Manfaat Akademik:}
    \begin{enumerate}
        \item Memberikan kontribusi terhadap pengembangan ilmu pengetahuan di bidang \textbf{kecerdasan buatan (Artificial Intelligence)} dan \textbf{pemrosesan suara (Speech Processing)}, khususnya dalam konteks \textit{Automatic Speech Recognition} (ASR) untuk bahasa dengan sumber daya terbatas.
        \item Menjadi \textbf{referensi ilmiah} dan bahan pembelajaran bagi mahasiswa atau peneliti lain yang tertarik untuk melakukan penelitian serupa dalam bidang \textit{speech-to-text}, \textit{low-resource language modeling}, dan \textit{model adaptation}.
        \item Menunjukkan penerapan teknik \textit{fine-tuning} pada model \textbf{OpenAI Whisper Small} menggunakan pustaka \textbf{PyTorch} dan \textbf{Hugging Face Transformers}, sehingga dapat menjadi dasar penelitian lanjutan terkait efisiensi model, optimasi pelatihan, dan adaptasi multibahasa.
    \end{enumerate}

    \item \textbf{Manfaat Praktis:}
    \begin{enumerate}
        \item Menghasilkan \textbf{sistem pengenalan suara bahasa Minangkabau} yang dapat digunakan sebagai alat bantu transkripsi otomatis untuk konten audio lokal, seperti wawancara, podcast, atau dokumentasi budaya.
        \item Mendukung \textbf{pelestarian bahasa daerah} melalui pemanfaatan teknologi AI yang dapat mempermudah digitalisasi dan dokumentasi bahasa Minangkabau.
        \item Memberikan \textbf{contoh implementasi praktis} penggunaan model besar berbasis \textit{Transformer} dalam konteks lokal, yang dapat menjadi acuan bagi pengembang, lembaga penelitian, dan institusi pendidikan di Indonesia.
    \end{enumerate}
\end{enumerate}


\section{Sistematika Penulisan} \label{I.Sistematika}
Sistematika penulisan berisi penjelasan mengenai susunan bab dalam laporan tugas akhir ini. Setiap bab disusun secara sistematis agar pembaca dapat memahami alur penelitian yang dilakukan, mulai dari tahap perumusan masalah hingga kesimpulan akhir. Adapun sistematika penulisan tugas akhir ini adalah sebagai berikut:

\subsection*{Bab I: Pendahuluan}
Bab ini berisi penjelasan mengenai \textbf{latar belakang penelitian}, \textbf{rumusan masalah}, \textbf{tujuan penelitian}, \textbf{batasan masalah}, \textbf{manfaat penelitian}, serta \textbf{sistematika penulisan}. Bab ini memberikan gambaran umum mengenai alasan dan arah penelitian yang dilakukan, sehingga pembaca memahami konteks dan ruang lingkup penelitian ini.

\subsection*{Bab II: Tinjauan Pustaka}
Bab ini membahas \textbf{landasan teori} dan \textbf{penelitian terdahulu} yang relevan dengan topik tugas akhir, meliputi konsep dasar \textit{Automatic Speech Recognition (ASR)}, model \textbf{OpenAI Whisper}, teknik \textit{fine-tuning} dan \textit{parameter-efficient fine-tuning} seperti \textbf{LoRA}, serta pembahasan mengenai \textbf{bahasa Minangkabau sebagai bahasa low-resource}. Bab ini menjadi dasar konseptual dan teoritis bagi metode yang digunakan dalam penelitian.

\subsection*{Bab III: Metodologi Penelitian}
Bab ini menjelaskan \textbf{tahapan penelitian} yang dilakukan, termasuk rancangan sistem, arsitektur model, dataset yang digunakan, tahapan \textit{fine-tuning}, lingkungan pengujian, serta metode evaluasi performa. Selain itu, bab ini juga menjelaskan alur kerja penelitian mulai dari pengumpulan data, pelatihan model, hingga proses evaluasi hasil.

\subsection*{Bab IV: Implementasi dan Hasil Penelitian}
Bab ini membahas \textbf{implementasi model dan sistem} yang dikembangkan, hasil pengujian terhadap model \textit{Whisper Small} yang telah di-\textit{fine-tune}, serta analisis performa berdasarkan metrik \textit{Word Error Rate (WER)}, \textit{Character Error Rate (CER)}, dan evaluasi kualitatif terhadap hasil transkripsi. Pada bagian ini juga dijelaskan perbandingan antara model dasar dan model hasil adaptasi terhadap bahasa Minangkabau.

\subsection*{Bab V: Kesimpulan dan Saran}
Bab ini berisi \textbf{kesimpulan} yang diperoleh dari hasil penelitian dan \textbf{saran} untuk pengembangan lebih lanjut. Kesimpulan mencakup hasil utama penelitian, tingkat keberhasilan implementasi sistem ASR bahasa Minangkabau, serta potensi pengembangan model atau sistem di masa depan.
